\documentclass[11pt]{article}
\usepackage[T1]{fontenc}
\usepackage[margin=1in]{geometry}
\usepackage{amsmath,amssymb,amsthm}
\usepackage{graphicx}
\usepackage{booktabs}
\usepackage[numbers]{natbib}
\usepackage{authblk}
\usepackage[colorlinks,allcolors=blue]{hyperref}

\newtheorem{theorem}{Theorem}
\newtheorem{lemma}[theorem]{Lemma}
\newtheorem{assumption}{Assumption}
\theoremstyle{definition}
\newtheorem{definition}[theorem]{Definition}
\theoremstyle{remark}
\newtheorem*{observation}{Observation}

\title{Title of the Paper}
\author[1]{First Author}
\author[2]{Second Author}
\affil[1]{Department, University, City, Country\\\texttt{first@example.org}}
\affil[2]{Company, City, Country\\\texttt{second@example.com}}
\date{}

\begin{document}
\maketitle

\begin{abstract}
State the problem, what you did, and the main result in a few sentences.
\end{abstract}

\noindent\textbf{Keywords:} first keyword, second keyword, third keyword

\section{Introduction}
\label{sec:intro}
Cite related work like this~\cite{knuth}, or several works at once~\cite{knuth,lamport}.

\begin{definition}
A thing is \emph{good} if it works.
\end{definition}

\begin{assumption}
Things work.
\end{assumption}

\begin{theorem}
\label{thm:main}
Every good thing works.
\end{theorem}

\begin{lemma}
Good things exist.
\end{lemma}

\begin{proof}[Proof of Theorem~\ref{thm:main}]
By definition,
\[
  1 + 1 = 2. \qedhere
\]
\end{proof}

\begin{observation}
This is plain.
\end{observation}

\begin{figure}[htbp]
  \centering
  \includegraphics[width=0.6\linewidth]{fig.png}
  \caption{Overview of the method.}
  \label{fig:overview}
\end{figure}

\section{Conclusion}
Summarise the findings and what comes next.

\section*{Acknowledgements}
Thank the people and funding that supported this work.

\bibliographystyle{plainnat}
\bibliography{refs}

\appendix
\section{Additional Results}
Material that supports the paper but is not needed to follow it.

\end{document}
